\chapter{Introducción y Fundamentación Teórica}

La instrumentación electrónica y los sistemas de adquisición de datos constituyen el núcleo de los procesos de control,
ensayo y automatización en la ingeniería contemporánea. En las aplicaciones de ensayo mecánico e industrial, las
magnitudes físicas primarias (tales como fuerza, desplazamiento, presión hidráulica y par torsional) no son eléctricas
en su origen. Por lo tanto, deben ser detectadas mediante sensores y transductores apropiados, acondicionadas
analógicamente para optimizar su rango dinámico y relación señal a ruido, y finalmente digitalizadas mediante sistemas
embebidos para su procesamiento y registro.

Uno de los mayores desafíos en el diseño e implementación de sistemas de instrumentación sensibles radica en la
convivencia con perturbaciones e interferencias electromagnéticas. Las señales provenientes de transductores
piezorresistivos (como las celdas de carga y puentes de galgas extensométricas) entregan tensiones en el orden de los
milivoltios o microvoltios. Si no se aplican criterios rigurosos de apantallamiento, blindaje y referencia a tierra,
las señales útiles quedan completamente degradadas o enmascaradas por el ruido ambiente y las interferencias de línea o
de maquinaria rotativa cercana.

\section{Objetivos del Trabajo Práctico}
\begin{itemize}
    \item Relevar in situ y diseñar la cadena de instrumentación y acondicionamiento para una máquina de ensayo de
    tracción y compresión, y para un banco dinamométrico de motores de combustión interna, a partir de sensores
    industriales reales.
    \item Diseñar las etapas analógicas de acondicionamiento de señal (adaptación de nivel, atenuación, filtrado
    pasabajos y amplificación de instrumentación) para adecuar los rangos de tensión a los conversores ADC.
    \item Seleccionar el sistema embebido para la digitalización y transmisión a PC, calculando la resolución mínima
    requerida en bits del conversor analógico-digital (ADC), la frecuencia de muestreo y la tensión de referencia.
    \item Cuantificar experimentalmente el acoplamiento de perturbaciones de baja frecuencia provenientes de la red
    eléctrica comercial ($\qty{220}{\volt}$ a $\qty{50}{\hertz}$).
    \item Evaluar el comportamiento estático de una celda de carga diferencial frente a diversas cargas y contrastar
    las limitaciones de la medición directa sin acondicionamiento frente a la medición acondicionada.
    \item Analizar experimentalmente el impacto de interferencias electromagnéticas severas generadas por un motor
    universal con escobillas sobre la etapa de amplificación.
    \item Evaluar cuantitativa y cualitativamente la efectividad de los métodos de mitigación: apantallamiento con
    gabinete metálico (jaula de Faraday) y conexión de la malla de cables blindados a masa frente a condiciones de
    malla flotante y circuito sin blindaje.
\end{itemize}

\section{Fundamentación Teórica}

\subsection{Distinción entre Ruido e Interferencia}
En los sistemas de medición electrónica es imprescindible diferenciar conceptualmente el ruido intrínseco de la
interferencia electromagnética externa:
\begin{enumerate}
    \item \emph{Ruido}: Corresponde a fluctuaciones eléctricas aleatorias originadas en procesos físicos microscópicos
    espontáneos dentro de los propios componentes del circuito. Entre sus manifestaciones primarias se encuentran el
    ruido térmico de Johnson-Nyquist generado por la agitación térmica de los portadores de carga en conductores y
    resistores, el ruido de disparo (\emph{shot noise}) debido a la naturaleza discreta de la carga al atravesar
    barreras de potencial en junturas semiconductoras, y el ruido de parpadeo (\emph{flicker} o $1/f$). El ruido es de
    naturaleza estocástica e impredecible; no puede eliminarse, únicamente minimizarse mediante la selección de
    componentes de bajo ruido, limitación del ancho de banda y reducción de la temperatura.
    \item \emph{Interferencia}: Comprende todas aquellas señales eléctricas espurias no deseadas que se acoplan al
    sistema de medición desde fuentes externas artificiales o naturales. Al estar originadas en fenómenos causales
    determinísticos (tales como la red de energía de $\qty{50}{\hertz}$, arcos eléctricos en colectores de delgas,
    conmutación en fuentes switching o radiotransmisiones), las interferencias son predecibles y factibles de ser
    suprimidas o atenuadas drásticamente mediante un diseño electromagnético adecuado.
\end{enumerate}

\subsection{Mecanismos de Acoplamiento de Perturbaciones}
Las interferencias externas ingresan a los circuitos electrónicos a través de cuatro mecanismos principales:
\begin{itemize}
    \item \emph{Acoplamiento capacitivo (campo eléctrico)}: Producido por la existencia de capacidades parásitas
    mutuas entre conductores a potenciales elevados y las líneas de señal. La corriente inyectada es proporcional a la
    capacidad parásita y a la derivada temporal de la tensión perturbadora ($i = C \frac{dv}{dt}$).
    \item \emph{Acoplamiento inductivo (campo magnético)}: Ocurre cuando un bucle formado por los conductores de señal
    concatena un flujo magnético variable en el tiempo generado por corrientes elevadas cercanas. De acuerdo con la ley
    de inducción de Faraday, se induce una fuerza electromotriz espuria en serie con el bucle:
    \begin{equation}
        e_{\text{ind}} = -\frac{d\Phi_B}{dt} = -M \frac{di}{dt}
        \label{eq.ej3:induccion_faraday}
    \end{equation}
    donde $M$ es la inductancia mutua y $\Phi_B$ es el flujo magnético enlazado por el área del circuito.
    \item \emph{Acoplamiento galvánico (por impedancia común)}: Ocurre cuando dos o más circuitos comparten un tramo
    común de retorno de corriente (conductor o pista de masa). La caída de tensión producida por la corriente de un
    circuito sobre la resistencia o inductancia parásita de masa se suma directamente a la señal del otro circuito.
    \item \emph{Acoplamiento por radiación electromagnética}: Se presenta a frecuencias medias y altas, donde las
    pistas del circuito o los cables de conexión poseen dimensiones comparables a una fracción de la longitud de onda
    ($\lambda/4$), comportándose como antenas receptoras.
\end{itemize}

\subsection{Celdas de Carga y Puente de Wheatstone}
Una celda de carga resistiva típica está compuesta por cuatro galgas extensométricas fijadas a un elemento elástico
deformable (viga de aluminio o acero) conectadas en configuración de puente de Wheatstone completo. Al aplicarse una
fuerza externa, dos galgas experimentan tracción ($+\Delta R$) y dos compresión ($-\Delta R$), manteniendo la simetría
y compensando automáticamente las variaciones térmicas.

Bajo una tensión de excitación continua $V_E$, la tensión diferencial en bornes del puente en condición balanceada
ideal viene dada por:
\begin{equation}
    V_d = V^+ - V^- = V_E \cdot \left( \frac{\Delta R}{R} \right)
    \label{eq.ej4:puente_balanceado}
\end{equation}
Dado que la deformación unitaria $\varepsilon = \frac{\Delta L}{L}$ en régimen elástico es diminuta, la variación
fraccional de resistencia $\frac{\Delta R}{R} = G_F \cdot \varepsilon$ suele ser del orden de $10^{-3}$, lo que
significa que para una excitación $V_E = \qty{5}{\volt}$, la señal máxima de salida oscila entre
$\qty{5}{\milli\volt}$ y $\qty{10}{\milli\volt}$.

\subsection{Amplificador de Instrumentación y Rechazo de Modo Común (CMRR)}
Para procesar señales diferenciales de baja amplitud inmersas en tensiones continuas de modo común elevadas
($V_{CM} \approx V_E / 2 = \qty{2.5}{\volt}$) y en presencia de ruido de línea acoplado en modo común, resulta
mandatorio el uso de un amplificador de instrumentación. Este circuito provee una amplificación diferencial
precisa $A_d$ y una atenuación severa de las señales presentes simultáneamente en ambas entradas ($A_{CM} \approx 0$).

La figura de mérito primordial es la Relación de Rechazo en Modo Común (CMRR), definida como:
\begin{equation}
    \text{CMRR} = 20 \log_{10}\left( \frac{|A_d|}{|A_{CM}|} \right)
    \label{eq.ej4:cmrr_def}
\end{equation}
Dispositivos integrados de alta precisión como el INA128 o INA122 ofrecen valores de CMRR superiores a
$100\si{\deci\bel}$, junto con impedancias de entrada en modo común y diferencial superiores a $10^{10}\ohm$,
previniendo cualquier alteración en las ramas resistivas del puente de Wheatstone. La ganancia se fija mediante un
único resistor externo $R_G$.

\subsection{Técnicas de Blindaje y Rectificación por Interferencia de Radiofrecuencia (RFI)}
El apantallamiento electrostático mediante chasis o gabinetes metálicos conductores confina el circuito en un volumen
equipotencial (caja de Faraday), bloqueando las líneas de campo eléctrico externo. Para que este blindaje cumpla su
función de derivación, debe estar vinculado al potencial de referencia del sistema (GND). Si el blindaje se deja
desconectado o flotante, no sólo pierde eficacia, sino que puede incrementar el ruido por acoplamiento capacitivo
cruzado.

Por otra parte, cuando trenes de impulsos de alta frecuencia producidos por descargas de conmutación (por ejemplo, el
chisporroteo entre escobillas y delgas de un motor universal) inciden sobre un amplificador operacional, se suscita el
fenómeno de \emph{demodulación o rectificación por EMI}. Las uniones colector-base o fuente-drenaje de los transistores
de entrada actúan como detectores de envolvente no lineales ante perturbaciones de radiofrecuencia que exceden el ancho
de banda lineal del circuito. Esta rectificación no lineal genera corrientes continuas espurias internas que alteran la
tensión de polarización, desplazando de forma sustancial el nivel de continua (\emph{DC offset shift}) en la salida.
